% !TEX root = ../journal_0421052099.tex
%
% Condensed from the results half of thesis_book/chapters/evaluation.tex. Every
% number is the book's, which reads it from results/phase4/thesis_numbers.json.
\section{Results}\label{sec:results}

\begin{table*}[!t]
\caption{Main results, $n = 400$ per dataset. Brackets are $95\%$ bootstrap
  intervals on Hits@1. P and R are mean per-question precision and recall, and
  Hedge is the share of questions with an empty answer. The Hits@1 ceilings are
  $97.0\%$ and $99.2\%$.}
\label{tab:main}
\centering
\small
\setlength{\tabcolsep}{4.5pt}
\begin{tabular}{@{}llrrrrlrrrr@{}}
\toprule
& \multicolumn{5}{c}{WebQSP} & \multicolumn{5}{c}{CWQ} \\
\cmidrule(lr){2-6} \cmidrule(l){7-11}
System & Hits@1 & F1 & P & R & Hedge & Hits@1 & F1 & P & R & Hedge \\
\midrule
No-retrieval & $0.453$ [$0.403$, $0.500$] & $0.305$ & $0.380$ & $0.303$ & $12.2\%$ & $0.307$ [$0.263$, $0.352$] & $0.279$ & $0.285$ & $0.282$ & $38.0\%$ \\
Vector-RAG & $0.568$ [$0.517$, $0.618$] & $0.432$ & $0.469$ & $0.460$ & $27.0\%$ & $0.203$ [$0.163$, $0.242$] & $0.168$ & $0.171$ & $0.189$ & $48.2\%$ \\
GraphRAG & $0.340$ [$0.292$, $0.388$] & $0.262$ & $0.279$ & $0.284$ & $55.8\%$ & $0.205$ [$0.168$, $0.245$] & $0.183$ & $0.187$ & $0.194$ & $60.8\%$ \\
Think-on-Graph & $0.660$ [$0.613$, $0.705$] & $0.477$ & $0.571$ & $0.480$ & $18.8\%$ & $0.435$ [$0.388$, $0.482$] & $0.378$ & $0.403$ & $0.384$ & $\mathbf{22.5\%}$ \\
AGR & $\mathbf{0.755}$ [$0.713$, $0.797$] & $\mathbf{0.642}$ & $\mathbf{0.697}$ & $\mathbf{0.653}$ & $\mathbf{8.2\%}$ & $\mathbf{0.522}$ [$0.475$, $0.570$] & $\mathbf{0.469}$ & $\mathbf{0.484}$ & $\mathbf{0.481}$ & $23.0\%$ \\
\bottomrule
\end{tabular}
\end{table*}

\subsection{Main Results}\label{sec:main}

On WebQSP, AGR answers $302$ of $400$ questions correctly, a Hits@1 of $0.755$
against a ceiling of $97.0\%$ (\Cref{tab:main}). That is $9.5$ points above
Think-on-Graph and $30.2$ above the parametric floor. The floor matters. The
backbone answers $45.3\%$ of these questions from memory alone. So any claim
about what retrieval contributes starts from $0.453$, not from zero.

F1 widens the margin that Hits@1 understates. Against Think-on-Graph the Hits@1
gap is $9.5$ points and the F1 gap $16.5$. The precision column, at $0.697$
against $0.571$, shows where the extra margin comes from. WebQSP questions carry
seven gold answers on average. So a system that asserts a long list collects
Hits@1 credit for the one entity that matches and pays nothing for the rest.
Precision charges for the rest. Together the two columns say that AGR answers
more questions and asserts fewer falsehoods. It does so at an $8.2\%$ hedge
rate, the lowest of the five systems, so its precision is not bought with
abstention.

Raw hits mislead on the retrieval baselines. GraphRAG scores $0.340$, below the
no-retrieval control's $0.453$. But the control asserts on $351$ questions and
is wrong on $170$, an assertion precision of $51.6\%$. GraphRAG asserts on $177$
and is wrong on $41$, for $76.8\%$. The retrieval baselines are told to answer
only from the facts given, so they abstain when retrieval misses. The control
guesses from memory, and on a memorised benchmark guessing pays. A grounded
system that scores below an ungrounded one on raw hits, while committing four
times fewer falsehoods, is the contamination the control was built to expose.

On CWQ, AGR answers $209$ of $400$, a Hits@1 of $0.522$ against a ceiling of
$99.2\%$. That is $8.7$ points above Think-on-Graph and $21.5$ above the floor.
Both static baselines now fall below the parametric control. Vector-RAG drops
from $0.568$ to $0.203$ and GraphRAG from $0.340$ to $0.205$, while no-retrieval
falls only from $0.453$ to $0.307$. Vector-RAG's context is one verbalised
triple, which cannot hold a chain at any radius. GraphRAG's one-hop radius
confounds its result, so the claim rests on Vector-RAG. The two agentic systems
commit at nearly the same rate, on $308$ and $310$ questions. From those
commitments AGR is right $209$ times and wrong $99$, against $174$ and $136$ for
Think-on-Graph. That is an assertion precision of $67.9\%$ against $56.1\%$.

All eight paired comparisons reject (\Cref{tab:mcnemar}). That includes the two
against Think-on-Graph, the only baseline that navigates the same graph with the
same model under the same budget. On CWQ, AGR uniquely solves $86$ questions
that Think-on-Graph misses, against $51$ the other way.

\begin{table}[!t]
\caption{Paired McNemar tests of AGR against each baseline on the same $400$
  questions. ``AGR only'' counts the questions that only AGR answers correctly,
  and ``Base only'' counts the reverse.}
\label{tab:mcnemar}
\centering
\small
\begin{tabular}{@{}lrrl@{}}
\toprule
Baseline & AGR only & Base only & $p$ \\
\midrule
\multicolumn{4}{@{}l}{\emph{WebQSP}} \\
No-retrieval & $152$ & $31$ & $2.3 \times 10^{-20}$ \\
Vector-RAG & $117$ & $42$ & $2.2 \times 10^{-9}$ \\
GraphRAG & $186$ & $20$ & $6.5 \times 10^{-35}$ \\
Think-on-Graph & $76$ & $38$ & $4.8 \times 10^{-4}$ \\
\midrule
\multicolumn{4}{@{}l}{\emph{CWQ}} \\
No-retrieval & $121$ & $35$ & $2.7 \times 10^{-12}$ \\
Vector-RAG & $151$ & $23$ & $2.8 \times 10^{-24}$ \\
GraphRAG & $148$ & $21$ & $1.0 \times 10^{-24}$ \\
Think-on-Graph & $86$ & $51$ & $3.5 \times 10^{-3}$ \\
\bottomrule
\end{tabular}
\end{table}

\subsection{Where the Margin over the Agentic Baseline Comes
From}\label{sec:margin}

The shared call cap binds on Think-on-Graph far more often than on AGR, because
beam search costs width times depth. So the same budget is not the same freedom
to search. Splitting the questions on whether the cap actually cut
Think-on-Graph off separates search quality from affordability
(\Cref{tab:split}).

\begin{table}[!t]
\caption{The comparison against Think-on-Graph, split on whether the shared
  $25$-call cap cut it short. The combined rows reproduce the Hits@1 column of
  \Cref{tab:main}.}
\label{tab:split}
\centering
\small
\begin{tabular}{@{}lrrrr@{}}
\toprule
Subset & $n$ & ToG & AGR & $\Delta$ \\
\midrule
\multicolumn{5}{@{}l}{\emph{WebQSP}} \\
ToG ran to completion & $283$ & $0.852$ & $0.788$ & $-0.064$ \\
ToG clipped by the cap & $117$ & $0.197$ & $0.675$ & $+0.478$ \\
Combined & $400$ & $0.660$ & $0.755$ & $+0.095$ \\
\midrule
\multicolumn{5}{@{}l}{\emph{CWQ}} \\
ToG ran to completion & $224$ & $0.629$ & $0.607$ & $-0.022$ \\
ToG clipped by the cap & $176$ & $0.188$ & $0.415$ & $+0.227$ \\
Combined & $400$ & $0.435$ & $0.522$ & $+0.087$ \\
\bottomrule
\end{tabular}
\end{table}

On the questions it is allowed to finish, Think-on-Graph is ahead on both
datasets. It scores $0.852$ against AGR's $0.788$ on WebQSP and $0.629$ against
$0.607$ on CWQ. The entire aggregate margin comes from the $29.2\%$ and $44.0\%$
of questions where the cap truncates it. There its accuracy falls to $0.197$ and
$0.188$, while AGR's holds at $0.675$ and $0.415$.

So the finding is not that guided best-first search resolves entities better
than model-pruned beam search. On this evidence it does not. The finding is that
beam search cannot afford to finish under a budget a deployed system would
impose, and guided search can. Width times depth multiplies, and every question
pays it. A plan plus a scored frontier does not multiply. The budget was fixed
in advance because cost is a property of the method, not a nuisance to control
away.

Two qualifications keep the claim honest. The cap of $25$ calls is itself a
choice, and a larger one would move the split. And Think-on-Graph prunes from
narrower candidate sets than AGR (\Cref{sec:baselines}). That cannot explain the
clipping, since a thinner candidate set makes each step cheaper. But it does
make the unclipped scores of $0.852$ and $0.629$ a lower bound on what the
baseline's search could resolve with AGR's widths.

The breadth of AGR's answers is checked rather than assumed. On answered WebQSP
questions AGR names $3.37$ entities against Think-on-Graph's $2.56$, and a
metric that scores a hit on any match rewards naming more. But AGR's entity
precision on the same questions is higher too, at $0.760$ against $0.703$. On
CWQ the figures are $1.54$ against $1.40$ entities and $0.629$ against $0.520$
precision. If AGR bought breadth without accuracy, its precision would fall as
its entity count rose. It does not fall.

\subsection{Positioning against Published RoG Results}\label{sec:rog}

The environment is built from the subgraphs that RoG~\cite{luo2024reasoning}
publishes, so the two systems invite a direct comparison. \Cref{tab:rog} gives
it, and it does not favour this work. RoG leads on all four figures, and the
margin is not sampling noise. AGR's $95\%$ intervals on Hits@1 are $[71.3,
79.7]$ and $[47.5, 57.0]$, and RoG's $85.7$ and $62.6$ fall outside both.

\begin{table}[!t]
\caption{AGR against RoG's published figures, in points. RoG's row is measured
  over the full test splits and AGR's over the certified samples. The last row
  scores AGR's answers with RoG's released scorer.}
\label{tab:rog}
\centering
\small
\begin{tabular}{@{}lrrrr@{}}
\toprule
& \multicolumn{2}{c}{WebQSP} & \multicolumn{2}{c}{CWQ} \\
\cmidrule(lr){2-3} \cmidrule(l){4-5}
System & Hits@1 & F1 & Hits@1 & F1 \\
\midrule
RoG (published) & $85.7$ & $70.8$ & $62.6$ & $56.2$ \\
AGR & $75.5$ & $64.2$ & $52.2$ & $46.9$ \\
AGR, scored by RoG's code & $76.0$ & $64.9$ & $52.2$ & $47.2$ \\
\bottomrule
\end{tabular}
\end{table}

Four differences bound what the comparison establishes. The first is decisive.
RoG is trained on these benchmarks, with its backbone fine-tuned on training
splits that hold $2{,}826$ and $27{,}639$ questions in this distribution. AGR
performs no training of any kind. RoG also searches each question's own
subgraph, where AGR searches the union of them all. The backbones differ. And
RoG reports over the full test splits rather than over samples of them. A fifth
difference, the scorer, is measured and small. RoG's released scorer counts a
gold answer found anywhere inside the predicted text. Scored with its functions,
no figure of AGR's moves by as much as a point (\Cref{tab:rog}). So the
comparison places AGR. A training-free system reaches $75.5$ and $52.2$ against
a supervised system's $85.7$ and $62.6$, on the same data, with no in-domain
examples. The contribution of this work lies outside the accuracy column.

\subsection{Accuracy by Hop Count}\label{sec:hops}

\begin{figure*}[!t]
\centering
% !TEX root = ../thesis_book_0421052099.tex
% GENERATED by scripts/build_figures.py from
% results/phase4/thesis_numbers.json -- do not edit.

\definecolor{agrNoRet}{HTML}{6E6E6E}
\definecolor{agrVector}{HTML}{E69F00}
\definecolor{agrGraph}{HTML}{CC79A7}
\definecolor{agrToG}{HTML}{009E73}
\definecolor{agrAGR}{HTML}{0072B2}
\definecolor{agrWQ}{HTML}{0072B2}
\definecolor{agrCWQ}{HTML}{D55E00}
\begin{tikzpicture}
\begin{axis}[agrplot, name=hop1, 
  width=0.46\textwidth, height=52mm,
  title={WebQSP},
  xlabel={Hop stratum}, ylabel={Hits@1},
  xtick={0,1,2}, xticklabels={{1 hop ($n{=}256$)}, {2 hops ($n{=}128$)},
    {3+ hops ($n{=}4$)}},
  xticklabel style={align=center},
  xmin=-0.25, xmax=2.25, ymin=0, ymax=0.9,
  legend style={at={(1.08,-0.30)}, anchor=north, legend columns=5,
    column sep=1.8mm},
  legend cell align=left,
  legend image post style={solid},
]
\addplot[dashed, mark=o, mark size=2.2pt, color=agrNoRet, very thick]
  coordinates {(0,0.43) (1,0.52) (2,0.25)};
\addlegendentry{No-retrieval}
\addplot[dashed, mark=square, mark size=2.2pt, color=agrVector, very thick]
  coordinates {(0,0.82) (1,0.12) (2,0.0)};
\addlegendentry{Vector-RAG}
\addplot[dashed, mark=triangle, mark size=2.2pt, color=agrGraph, very thick]
  coordinates {(0,0.3) (1,0.44) (2,0.25)};
\addlegendentry{GraphRAG}
\addplot[dashed, mark=diamond, mark size=2.2pt, color=agrToG, very thick]
  coordinates {(0,0.63) (1,0.78) (2,0.25)};
\addlegendentry{Think-on-Graph}
\addplot[dashed, mark=*, mark size=2.2pt, color=agrAGR, very thick]
  coordinates {(0,0.75) (1,0.84) (2,0.0)};
\addlegendentry{AGR}
\end{axis}
\begin{axis}[agrplot, name=hop2, at={($(hop1.east)+(12mm,0)$)}, anchor=west,
  width=0.46\textwidth, height=52mm,
  title={ComplexWebQuestions},
  xlabel={Hop stratum}, ylabel={Hits@1},
  xtick={0,1,2}, xticklabels={{1 hop ($n{=}137$)}, {2 hops ($n{=}211$)},
    {3+ hops ($n{=}49$)}},
  xticklabel style={align=center},
  xmin=-0.25, xmax=2.25, ymin=0, ymax=0.9,
%
]
\addplot[mark=o, mark size=2.2pt, color=agrNoRet, very thick]
  coordinates {(0,0.38) (1,0.3) (2,0.14)};
\addplot[mark=square, mark size=2.2pt, color=agrVector, very thick]
  coordinates {(0,0.33) (1,0.14) (2,0.1)};
\addplot[mark=triangle, mark size=2.2pt, color=agrGraph, very thick]
  coordinates {(0,0.34) (1,0.16) (2,0.02)};
\addplot[mark=diamond, mark size=2.2pt, color=agrToG, very thick]
  coordinates {(0,0.53) (1,0.37) (2,0.45)};
\addplot[mark=*, mark size=2.2pt, color=agrAGR, very thick]
  coordinates {(0,0.46) (1,0.55) (2,0.57)};
\end{axis}
\end{tikzpicture}

\caption{Hits@1 across hop strata. On ComplexWebQuestions, AGR is the only
  system whose accuracy rises with hop count. WebQSP is drawn dashed because its
  stratum of three or more hops holds four questions and cannot carry a trend.
  The unreachable stratum is omitted, since it measures assertion rather than
  accuracy.}
\label{fig:hops}
\end{figure*}

\begin{table}[!t]
\caption{Hits@1 and F1 per hop stratum, written as Hits@1/F1. Stratum sizes are
  those of \Cref{tab:testsets}. WebQSP's stratum of three or more hops holds
  four questions and is not interpreted.}
\label{tab:strata}
\centering
\small
\setlength{\tabcolsep}{3pt}
\begin{tabular}{@{}lllll@{}}
\toprule
System & 1 hop & 2 hops & 3+ hops & Unreach. \\
\midrule
\multicolumn{5}{@{}l}{\emph{WebQSP}} \\
No-retrieval & $0.43$/$0.29$ & $0.52$/$0.36$ & $0.25$/$0.00$ & $0.17$/$0.17$ \\
Vector-RAG & $0.82$/$0.63$ & $0.12$/$0.07$ & $0.00$/$0.00$ & $0.08$/$0.08$ \\
GraphRAG & $0.30$/$0.23$ & $0.44$/$0.35$ & $0.25$/$0.07$ & $0.08$/$0.08$ \\
Think-on-Graph & $0.63$/$0.45$ & $0.78$/$0.58$ & $0.25$/$0.07$ & $0.17$/$0.17$ \\
AGR & $0.75$/$0.64$ & $\mathbf{0.84}$/$\mathbf{0.73}$ & $0.00$/$0.00$ & $0.08$/$0.08$ \\
\midrule
\multicolumn{5}{@{}l}{\emph{CWQ}} \\
No-retrieval & $0.38$/$0.35$ & $0.30$/$0.27$ & $0.14$/$0.14$ & $0.00$/$0.00$ \\
Vector-RAG & $0.33$/$0.27$ & $0.14$/$0.11$ & $0.10$/$0.09$ & $0.67$/$0.67$ \\
GraphRAG & $0.34$/$0.30$ & $0.16$/$0.14$ & $0.02$/$0.02$ & $0.33$/$0.33$ \\
Think-on-Graph & $\mathbf{0.53}$/$0.48$ & $0.37$/$0.32$ & $0.45$/$0.32$ & $0.33$/$0.33$ \\
AGR & $0.46$/$0.42$ & $\mathbf{0.55}$/$\mathbf{0.49}$ & $\mathbf{0.57}$/$\mathbf{0.50}$ & $0.33$/$0.33$ \\
\bottomrule
\end{tabular}
\end{table}

\Cref{fig:hops} and \Cref{tab:strata} answer RQ1 with a shape rather than a
difference of means. On CWQ, AGR's accuracy rises with hop count. Its Hits@1
goes $0.46 \to 0.55 \to 0.57$ across one, two, and three or more hops, and its
F1 goes $0.42 \to 0.49 \to 0.50$. It is the only system on that dataset that
ends above where it started. Three of the others decay monotonically.
No-retrieval goes $0.38 \to 0.30 \to 0.14$, Vector-RAG $0.33 \to 0.14 \to 0.10$,
and GraphRAG $0.34 \to 0.16 \to 0.02$. Think-on-Graph falls from $0.53$ to
$0.37$ and partly recovers to $0.45$, still below its own one-hop score. A
system that gets better as the required chain gets longer is not merely ahead on
average. Its profile is the only one that fits a system which actually traverses
the chain.

Two cells cut the other way, and they are reported for that reason.
Think-on-Graph beats AGR on CWQ's one-hop stratum, at $0.53$ against $0.46$ over
$137$ questions. The ablation of \Cref{sec:ablations} identifies the mechanism.
Decomposition costs accuracy on questions that do not need decomposing. And
Vector-RAG scores $0.82$ on WebQSP's one-hop stratum, the best single-hop number
anywhere, but $0.12$ at two hops. When one triple suffices, dense retrieval over
verbalised triples is excellent. When a chain is required, it has nothing to
offer.

GraphRAG's WebQSP profile is inverted, at $0.30$ on one hop against $0.44$ on
two. Part of that is hub noise. A popular entity's neighbourhood drowns the
answer, or the unordered sample of its edges leaves the answer out. The rest
marks the radius. WebQSP's two-hop stratum is largely mediator paths, which a
one-hop expansion reaches. CWQ's two-hop stratum is genuine composition, where
GraphRAG scores $0.16$. So its row measures what its radius covers, not how far
static retrieval can see.

The unreachable stratum is a hallucination probe rather than an accuracy
measurement. For these questions no path to gold exists within the hop bound. So
every hit there is an assertion the environment cannot support, one that
happened to match the label. Vector-RAG scores $0.67$ on CWQ's three unreachable
questions, and no-retrieval $0.17$ on WebQSP's twelve. The counts are tiny and
carry no weight. But the direction is the one \Cref{sec:groundedness} measures
at scale.

\subsection{Groundedness}\label{sec:groundedness}

\Cref{tab:groundedness} reports both tiers. The structural differential is
emphatic. The parametric control asserts $661$ entities on WebQSP, of which
$179$ ($27.1\%$) have no basis in the environment. Both navigating systems sit
at exactly zero, over $1{,}235$ entities for AGR and $832$ for Think-on-Graph.
Graph-navigated answering removes structural hallucination outright.

\begin{table}[!t]
\caption{Two-tier groundedness. Tier~1 counts asserted entities that are absent
  from the graph or more than four hops from every topic entity. Tier~2 is the
  judged share of $60$ sampled answers per cell that are supported as the
  answer.}
\label{tab:groundedness}
\centering
\small
\begin{tabular}{@{}lrr@{}}
\toprule
System & Tier 1 ungrounded & Tier 2 supported \\
\midrule
\multicolumn{3}{@{}l}{\emph{WebQSP}} \\
No-retrieval & $179$ of $661$ ($27.1\%$) & $63.3\%$ \\
Vector-RAG & $38$ of $1{,}040$ ($3.7\%$) & $61.7\%$ \\
GraphRAG & $6$ of $800$ ($0.8\%$) & $55.0\%$ \\
Think-on-Graph & $0$ of $832$ ($0.0\%$) & $61.7\%$ \\
AGR & $0$ of $1{,}235$ ($0.0\%$) & $66.7\%$ \\
\midrule
\multicolumn{3}{@{}l}{\emph{CWQ}} \\
No-retrieval & $42$ of $340$ ($12.4\%$) & $36.7\%$ \\
Vector-RAG & $7$ of $289$ ($2.4\%$) & $50.0\%$ \\
GraphRAG & $2$ of $224$ ($0.9\%$) & $36.7\%$ \\
Think-on-Graph & $0$ of $433$ ($0.0\%$) & $43.3\%$ \\
AGR & $0$ of $474$ ($0.0\%$) & $48.3\%$ \\
\bottomrule
\end{tabular}
\end{table}

But Think-on-Graph reaches zero as well, and that changes the attribution.
Structural groundedness is a property of the navigation paradigm. Any system
that copies its answer entities out of triples it traversed inherits it. It is
not a property of the verification layer. The two navigating systems differ at
the level of semantic support. Think-on-Graph's precision of $0.571$ and
$0.403$, against AGR's $0.697$ and $0.484$, shows it asserting more entities
that exist and connect to the topic but are not the answer. Vector-RAG's few
ungrounded entities are mostly retrieval errors rather than fabrications. Its
index is global, so retrieval can surface a similarly named subject elsewhere in
the graph. And the parametric control's CWQ rate of $12.4\%$ is less than half
its WebQSP rate. That is not because it fabricates less on harder questions. It
hedges more, at $38.0\%$ against $12.2\%$, and abstention pushes the measured
rate down.

Tier~2 places every system in a band from $36.7\%$ to $66.7\%$. AGR is highest
on WebQSP and second on CWQ, $5.0$ points ahead of Think-on-Graph on both.
Neither gap is resolvable at this sample size. With $60$ judgements per cell,
the $95\%$ interval has a half-width of about $12$ to $13$ points. What the tier
does establish is the level. Zero structural hallucination coexists with a third
to a half of assertions failing a strict test of semantic support. That is a
statement about the paradigm, not about any one system.

\subsection{Efficiency}\label{sec:efficiency}

Agency costs roughly ten times what one-shot retrieval costs (\Cref{tab:cost}).
AGR takes $4{,}511$ and $6{,}818$ tokens per question against about $500$, and
$11$ to $17$ seconds against one or two. That is the price of RQ1's accuracy,
and it is stated as a price.

\begin{table}[!t]
\caption{Cost per question. Seconds use cold-cache records only. Clipped is the
  share of questions on which the shared $25$-call budget ran out mid-search.}
\label{tab:cost}
\centering
\small
\begin{tabular}{@{}lrrrr@{}}
\toprule
System & Tokens & Calls & Seconds & Clipped \\
\midrule
\multicolumn{5}{@{}l}{\emph{WebQSP}} \\
No-retrieval & $113$ & $1.0$ & $1.2$ & --- \\
Vector-RAG & $527$ & $1.0$ & $1.4$ & --- \\
GraphRAG & $531$ & $1.0$ & $2.1$ & --- \\
Think-on-Graph & $3{,}615$ & $12.8$ & $12.9$ & $29.2\%$ \\
AGR & $4{,}511$ & $6.2$ & $11.2$ & $0.0\%$ \\
\midrule
\multicolumn{5}{@{}l}{\emph{CWQ}} \\
No-retrieval & $118$ & $1.0$ & $1.0$ & --- \\
Vector-RAG & $535$ & $1.0$ & $1.4$ & --- \\
GraphRAG & $470$ & $1.0$ & $2.2$ & --- \\
Think-on-Graph & $5{,}572$ & $18.2$ & $15.2$ & $44.0\%$ \\
AGR & $6{,}818$ & $8.9$ & $17.2$ & $0.0\%$ \\
\bottomrule
\end{tabular}
\end{table}

The informative comparison is with the other agentic system. AGR uses half the
calls, $6.2$ against $12.8$ on WebQSP and $8.9$ against $18.2$ on CWQ. It spends
more tokens per call, because each call carries a plan, a scored candidate set,
and a fact window. Under the shared budget, Think-on-Graph exhausts its calls on
$29.2\%$ and $44.0\%$ of questions, while AGR never reaches the call cap. AGR's
own budgets do bind. The depth cap refuses a deeper expansion on $25.6\%$ of
test questions, against $6.2\%$ for backtracks and $1.9\%$ for repair.

\begin{figure*}[!t]
\centering
% !TEX root = ../thesis_book_0421052099.tex
% GENERATED by scripts/build_figures.py from
% results/phase4/thesis_numbers.json -- do not edit.

\definecolor{agrNoRet}{HTML}{6E6E6E}
\definecolor{agrVector}{HTML}{E69F00}
\definecolor{agrGraph}{HTML}{CC79A7}
\definecolor{agrToG}{HTML}{009E73}
\definecolor{agrAGR}{HTML}{0072B2}
\definecolor{agrWQ}{HTML}{0072B2}
\definecolor{agrCWQ}{HTML}{D55E00}
\begin{tikzpicture}
\begin{axis}[agrplot, name=plot1, 
  width=0.46\textwidth, height=52mm,
  title={WebQSP},
  xmode=log, log basis x=10,
  xlabel={Mean tokens per question}, ylabel={Hits@1},
  xmin=80, xmax=12000, ymin=0.15, ymax=0.85,
  legend style={at={(1.08,-0.30)}, anchor=north, legend columns=5,
    column sep=1.8mm},
  legend cell align=left,
  legend image post style={solid},
]
\addplot[only marks, mark=o, mark size=2.6pt, color=agrNoRet, very thick]
  coordinates {(113,0.453)};
\addlegendentry{No-retrieval}
\addplot[only marks, mark=square, mark size=2.6pt, color=agrVector, very thick]
  coordinates {(527,0.568)};
\addlegendentry{Vector-RAG}
\addplot[only marks, mark=triangle, mark size=2.6pt, color=agrGraph, very thick]
  coordinates {(531,0.34)};
\addlegendentry{GraphRAG}
\addplot[only marks, mark=diamond, mark size=2.6pt, color=agrToG, very thick]
  coordinates {(3615,0.66)};
\addlegendentry{Think-on-Graph}
\addplot[only marks, mark=*, mark size=2.6pt, color=agrAGR, very thick]
  coordinates {(4511,0.755)};
\addlegendentry{AGR}
\end{axis}
\begin{axis}[agrplot, name=plot2, at={($(plot1.east)+(12mm,0)$)}, anchor=west,
  width=0.46\textwidth, height=52mm,
  title={ComplexWebQuestions},
  xmode=log, log basis x=10,
  xlabel={Mean tokens per question}, ylabel={Hits@1},
  xmin=80, xmax=12000, ymin=0.15, ymax=0.85,
%
]
\addplot[only marks, mark=o, mark size=2.6pt, color=agrNoRet, very thick]
  coordinates {(118,0.307)};
\addplot[only marks, mark=square, mark size=2.6pt, color=agrVector, very thick]
  coordinates {(535,0.203)};
\addplot[only marks, mark=triangle, mark size=2.6pt, color=agrGraph, very thick]
  coordinates {(470,0.205)};
\addplot[only marks, mark=diamond, mark size=2.6pt, color=agrToG, very thick]
  coordinates {(5572,0.435)};
\addplot[only marks, mark=*, mark size=2.6pt, color=agrAGR, very thick]
  coordinates {(6818,0.522)};
\end{axis}
\end{tikzpicture}

\caption{Accuracy against token cost, one point per system, on a logarithmic
  token axis. Think-on-Graph sits below and to the left of AGR, so it is a
  frontier point on tokens, not a dominated one. The domination claim holds on
  calls, which are not plotted and are given in \Cref{tab:cost}.}
\label{fig:cost}
\end{figure*}

The cost frontier depends on which cost is charged (\Cref{fig:cost}). On tokens,
Think-on-Graph is not dominated. It buys its lower accuracy at a lower token
price on both datasets. On calls it is dominated, with lower Hits@1 at about
twice the calls. Calls are what the budget meters, and it is the call budget
that clips the baseline. Vector-RAG on WebQSP is the cheap corner, the right
choice for shallow questions and useless for deep ones.

\subsection{Ablations}\label{sec:ablations}

\begin{table*}[!t]
\caption{Ablation conditions on stratified halves ($n = 200$ WebQSP, $n = 198$
  CWQ). The reference row is the frozen system restricted to the same questions.
  P is mean per-question precision, and $p$ is the exact McNemar test against
  the reference on per-question correctness.}
\label{tab:ablation}
\centering
\small
\setlength{\tabcolsep}{4pt}
\begin{tabular}{@{}lrrrrrrlrrrrrrl@{}}
\toprule
& \multicolumn{7}{c}{WebQSP} & \multicolumn{7}{c}{CWQ} \\
\cmidrule(lr){2-8} \cmidrule(l){9-15}
Condition & Hits@1 & F1 & P & Hedge & Tokens & Calls & $p$ & Hits@1 & F1 & P & Hedge & Tokens & Calls & $p$ \\
\midrule
Reference & $0.755$ & $0.659$ & $0.710$ & $8.5\%$ & $4{,}562$ & $6.2$ & --- & $0.495$ & $0.445$ & $0.461$ & $23.2\%$ & $7{,}105$ & $9.1$ & --- \\
No-planner & $\mathbf{0.830}$ & $\mathbf{0.742}$ & $0.788$ & $5.5\%$ & $3{,}159$ & $4.0$ & $0.006$ & $0.434$ & $0.398$ & $0.407$ & $26.8\%$ & $5{,}617$ & $6.7$ & $0.088$ \\
No-backtrack & $0.745$ & $0.646$ & $0.704$ & $9.5\%$ & $4{,}261$ & $5.8$ & $0.73$ & $0.490$ & $0.445$ & $0.456$ & $28.3\%$ & $5{,}792$ & $7.4$ & $1.00$ \\
No-verifier & $0.760$ & $0.663$ & $0.711$ & $8.0\%$ & $4{,}460$ & $6.1$ & $1.00$ & $0.490$ & $0.436$ & $0.451$ & $20.2\%$ & $6{,}704$ & $8.7$ & $1.00$ \\
Embedding-only & $0.740$ & $0.643$ & $0.701$ & $13.5\%$ & $3{,}413$ & $4.4$ & $0.66$ & $0.475$ & $0.437$ & $0.452$ & $28.3\%$ & $4{,}925$ & $5.9$ & $0.48$ \\
\bottomrule
\end{tabular}
\end{table*}

\Cref{tab:ablation} gives the four conditions. One reaches significance, and it
points the opposite way from what the design intended.

\textbf{The planner.} On WebQSP, removing the planner improves every axis at
once. Hits@1 goes from $0.755$ to $0.830$ and F1 from $0.659$ to $0.742$, while
tokens fall from $4{,}562$ to $3{,}159$ and calls from $6.2$ to $4.0$. McNemar
gives $p = 0.006$ on $6$ against $21$ discordant pairs. Without the planner, the
one-hop stratum rises from $0.75$ to $0.85$ and the two-hop stratum from $0.84$
to $0.89$. On CWQ the sign reverses, at $0.495 \to 0.434$ on $27$ against $15$
discordant pairs ($p = 0.088$), short of significance. There the one-hop stratum
still improves ($0.49 \to 0.54$), while the two-hop stratum collapses ($0.48 \to
0.34$). That is the pattern of decomposition that helps genuine composition and
hurts questions that do not need it. So the planner's value depends on question
structure. Applying it to every question regardless is the defect.

\textbf{Backtracking.} Hits@1 moves by $-0.010$ and $-0.005$, with $p = 0.73$
and $1.00$. The mechanism is not idle. The full runs record $63$ backtracks on
WebQSP and $199$ on CWQ. But this experiment cannot measure their accuracy
consequence. We report no detectable effect at $n \approx 200$, not a confirmed
null.

\textbf{Verification.} This is the cleanest null in the table. Across $398$
paired questions, removing claim verification changed the correctness of exactly
two. Hits@1 moves by $+0.005$ and $-0.005$, and precision by $+0.001$ and
$-0.010$. The design predicts as much. On $77$ of $80$ development questions the
layer certified the draft and left it alone, and a mechanism that mostly passes
cannot move an accuracy metric. Claim verification does not measurably change
how often AGR is right. It changes how often AGR is right for a reason it can
show.

\textbf{Embedding-only scoring.} Removing the model term costs $0.015$ and
$0.020$ Hits@1. The effect is negative on both datasets and short of
significance ($p = 0.66$ and $0.48$). The condition is also cheaper on both, so
the model term buys accuracy with tokens as designed. The development sweep,
where $\alpha = 1$ lost five hits on $80$ questions, remains the main evidence
for $\alpha = 0.7$. The condition also carries the most backtracks, $57$ and
$143$ against the reference's $38$ and $108$. Part of that rise is an
implementation artefact. With the model term disabled, the signal compared
against $\tau$ is taken over fewer candidates, so it falls after a ban where the
reference's holds steady. So the backtrack count is corroborating, not
load-bearing.

One condition of four reaching significance is the correct outcome to report. No
correction for multiple comparisons was pre-specified, so the $p$-values are
read as descriptive. Applied after the fact, and labelled as such, a Bonferroni
correction changes no claim. The planner's $p = 0.00592$ clears $0.05/8 =
0.00625$ within the eight ablation tests, and the eight baseline comparisons
clear the same bar, the largest of them at $p = 0.0035$. The power limit is
arithmetic. At $n \approx 200$ and accuracy between $0.5$ and $0.75$, a
component that changes five questions in a hundred produces on the order of ten
discordant pairs. That is not enough to reject. The planner cleared the bar
because its effect is large, at $21$ against $6$. So three conditions are
reported as having no detectable effect at this sample size, which is not the
same as no effect.

\endinput
